\documentclass[aspectratio=169]{beamer}
\usepackage{fontspec}

% Tipografías de la Identidad Visual
\setmainfont{Inter}
\setsansfont{Inter}
\setmonofont{JetBrains Mono}
\newfontfamily\titlefont{Montserrat}[BoldFont={Montserrat Bold}]

% Colores de la Identidad Visual
\definecolor{AzulMarino}{HTML}{1A365D}
\definecolor{GrisPizarra}{HTML}{4A5568}
\definecolor{BlancoNieve}{HTML}{F7FAFC}
\definecolor{VerdeAzulado}{HTML}{319795}

% Aplicación de colores a Beamer
\setbeamercolor{palette primary}{bg=AzulMarino, fg=BlancoNieve}
\setbeamercolor{palette secondary}{bg=VerdeAzulado, fg=BlancoNieve}
\setbeamercolor{palette tertiary}{bg=GrisPizarra, fg=BlancoNieve}
\setbeamercolor{structure}{fg=AzulMarino}
\setbeamercolor{background canvas}{bg=BlancoNieve}
\setbeamercolor{normal text}{fg=GrisPizarra}

\setbeamerfont{title}{family=\titlefont, series=\bfseries, size=\Huge}
\setbeamerfont{frametitle}{family=\titlefont, series=\bfseries, size=\LARGE}

\usepackage[utf8]{inputenc}
\usepackage{graphicx}
\usepackage{listings}

\lstset{
  basicstyle=\ttfamily\color{GrisPizarra}\small,
  keywordstyle=\color{VerdeAzulado}\bfseries,
  stringstyle=\color{AzulMarino},
  commentstyle=\color{GrisPizarra}\itshape,
  backgroundcolor=\color{BlancoNieve},
  frame=single,
  rulecolor=\color{GrisPizarra}
}

\title{Práctica Guiada: Limpieza y Transformación}
\subtitle{Scikit-Learn Preprocessing (Laboratorio)}
\author{Clase Práctica - Semana 2}
\date{}

\begin{document}

\begin{frame}
    \titlepage
\end{frame}

\begin{frame}[fragile]{Imputando Valores Nulos}
    Cuando los sensores fallan o los logs están corruptos:
    
\begin{lstlisting}[language=Python]
import pandas as pd

# 1. Calculamos la mediana segura
mediana = df['bytes_transmitted'].median()

# 2. Rellenamos (Fill NaN)
df['bytes_transmitted'] = df['bytes_transmitted'].fillna(mediana)

# 3. Borramos los registros donde falta la etiqueta objetivo
df = df.dropna(subset=['alert_type'])
\end{lstlisting}
    
    \begin{itemize}
        \item \textbf{\textcolor{VerdeAzulado}{fillna()}}: Asigna un valor a todo hueco vacío.
        \item \textbf{\textcolor{AzulMarino}{dropna()}}: Corta problemas de raíz eliminando filas.
    \end{itemize}
\end{frame}

\begin{frame}[fragile]{One-Hot Encoding Dinámico}
    Transformando categorías alfabéticas a la matrix binaria:
    
\begin{lstlisting}[language=Python]
df_encoded = pd.get_dummies(df, 
                            columns=['alert_type'], 
                            prefix='type', 
                            dtype=int)
\end{lstlisting}
    
    \begin{itemize}
        \item Crea columnas como \texttt{type\_Malware} con \texttt{1} o \texttt{0}.
        \item Listo para ser interpretado matemáticamente.
    \end{itemize}
\end{frame}

\begin{frame}[fragile]{Escaladores (StandardScaler)}
    Emparejando la cancha para el algoritmo k-NN:
    
\begin{lstlisting}[language=Python]
from sklearn.preprocessing import StandardScaler

scaler = StandardScaler()
# El algoritmo aprende (fit) y transforma la data
X_scaled = scaler.fit_transform(X)
\end{lstlisting}
    
    \begin{alertblock}{Regla de Oro en Machine Learning}
    ¡NUNCA apliques \texttt{fit} en tus datos de Prueba (\texttt{X\_test})! Solo usa \texttt{transform()}, de lo contrario se filtra información del futuro (Data Leakage).
    \end{alertblock}
\end{frame}

\end{document}
